\subsection{Shrinking boundary layers and the attempted-bit logarithm}
The fixed-layer bisection of Theorem~\ref{thm:rare-stationary} uses the
final probability accuracy at every level. It remains a valid controller.
The following additional result uses a coarser layer while the radial
interval is wide. A safeguarded interval update permits the layer to
shrink without losing the true boundary from the interval.

\begin{proposition}[A shrinking-layer rare-collision controller]
\label{prop:shrinking-rare}
Under the hypotheses of Theorem~\ref{thm:rare-stationary}, its geometric
and primitive-data conclusions can be achieved with
\begin{align}
 J_\nu&\le C\nu^{-1/(s-2)}\log(C/\nu),
                                  \label{eq:shrinking-centers}\\
 N_\nu&\le C\nu^{-((\gamma+3/2)s+1)/(s-2)}
                                      \log(C/(\nu\delta)).
                                  \label{eq:shrinking-attempts}
\end{align}
The assertion includes both the calibrated one-scale experiment and
the two known scales with unknown footprint, with their respective
fixed compass-separation hypotheses. Nominal centers require only the
same finest precision $O(\nu^{s/(s-2)})$. The launch spread, direction
law and fixed aperture are unchanged.
\end{proposition}
\begin{proof}
Reserve failure probability $\delta/2$ for the fixed-accuracy acquisition
of Lemma~\ref{lem:rare-coarse-normals}. We first describe one radial
search on this successful coarse event. Translate its interior center
to the origin, and write
\[
 \overline B_r\subset P\subset\overline B_{R_1},
 \qquad R=R_P(u)
\]
for its known uniform radial bounds and unknown boundary radius in the
chosen unit direction $u$. Fix a dyadic constant $A\ge4R_1/r$.
For all sufficiently small $e>0$, the inball inclusions
\begin{equation}\label{eq:shrinking-radial-conversion}
 (1-\rho/r)P\subset P_{-\rho},\qquad
 P+\rho\overline B\subset(1+\rho/r)P
                      \quad(0<\rho<r)
\end{equation}
show the following. If a nominal point within distance $e/2$ of $au$
is used in Proposition~\ref{prop:rare-query} with physical layer $e$,
its label is correct whenever $|a-R|>Ae$, except on the conditional
failure event of that proposition. Indeed $a<R-Ae$ places $au$ in
$P_{-2e}$ by the first inclusion with $\rho=2e$, whereas
$a>R+Ae$ places it outside $P+2e\overline B$ by the second.
The rounding reserve preserves depth at least $e$ in the first case
and an exterior point in the second. Thus, on a successful test,
\begin{equation}\label{eq:shrinking-label-invariant}
 \text{label }1\ \Longrightarrow\ a\le R+Ae,
 \qquad
 \text{label }0\ \Longrightarrow\ a\ge R-Ae.
\end{equation}
The fixed tubular neighborhood of the original coarse bracket supplies
the same valid candidate normal for these rounded points.

Take an initial certified interval $[l_0,u_0]$ containing $R$, with
fixed width $w_0>0$. Coarse accuracy and numerical reserves may be
chosen so that this width is a common dyadic prior constant for all
radial searches, and so that $w_0/(4A)$ lies below all collar and
mass-lemma thresholds. Suppose that a current interval $[l_j,u_j]$
contains $R$ and has width $w_j$. At its midpoint $a_j$ query with
\begin{equation}\label{eq:shrinking-layer-schedule}
            e_j=\frac{w_j}{4A},\qquad q=\frac34.
\end{equation}
For label one retain $[a_j-Ae_j,u_j]$; for label zero retain
$[l_j,a_j+Ae_j]$. Equation~\eqref{eq:shrinking-label-invariant}
proves that the retained interval still contains $R$. It is a subinterval
of the preceding interval and has width exactly $q w_j$, whichever
label occurs. Consequently
\[
                     w_j=w_0q^j.
\]
This deterministic width schedule holds also on an erroneous record;
only the containment assertion uses successful labels.

Put $h\asymp\nu^{1/(s-2)}$ and choose a positive dyadic final radial
tolerance $E\asymp c h^s$. Let $K$ be the least integer with
$w_K\le2E$.
For all sufficiently small $E$, one has $K\ge1$ and
\begin{equation}\label{eq:shrinking-final-width}
                   2qE<w_K\le2E,
             \qquad K\le C\log(C/E).
\end{equation}
The final midpoint has radial error at most $E$. Let $M_h\le C h^{-1}$
be a deterministic bound on the total number of radial searches,
including both scales when present. At level $j=0,\ldots,K-1$ of
each search use conditional failure allowance
\begin{equation}\label{eq:shrinking-confidence}
 \varepsilon_j=\frac{\delta}{4M_h}\,2^{-(K-1-j)}.
\end{equation}
The sum for one search is less than $\delta/(2M_h)$, so all fine
tests succeed, conditional on coarse success, with probability at least
$1-\delta/2$. This allocation permits relatively larger failure
allowances at the more expensive fine levels. Locations, labels and
the resulting subintervals may depend on the full past; fresh
preparations give the required conditional bounds at every level.

Write $\kappa=\gamma+3/2$. By Proposition~\ref{prop:rare-query},
the attempts in one search are at most
\begin{align*}
 C\sum_{j=0}^{K-1} e_j^{-\kappa}
                                  \log(C/\varepsilon_j)
 &\le C\sum_{j=0}^{K-1} w_j^{-\kappa}
          \left\{\log(CM_h/\delta)+(K-1-j)\log2\right\}\\
 &\le C w_K^{-\kappa}\log(CM_h/\delta).
\end{align*}
For the last inequality put $k=K-j$ and use the convergent sums
$\sum_{k\ge1}q^{\kappa k}$ and
$\sum_{k\ge1}kq^{\kappa k}$. Equation~\eqref{eq:shrinking-final-width}
therefore bounds the total fine cost by
\[
                         C M_h E^{-\kappa}\log(CM_h/\delta).
\]
The coarse cost is $C\log(C/\delta)$ and is absorbed by this bound.
Substitution of $M_h\le C h^{-1}$ and $E\asymp c h^s$ proves
\eqref{eq:shrinking-attempts}. There are at most two centers per
level and $M_hK$ levels in all; the coarse stencil count is fixed.
This proves \eqref{eq:shrinking-centers}.

All bracket endpoints, layer widths and updates are dyadic when $A$
is chosen as a power of two. For an explicit implementation, choose
one dyadic target mesh so that each nominal point is computed within
$E/(8A)$ of its ideal radial point. Since
$w_{K-1}>2E$, this error is less than $e_j/2$ at every level.
For sufficiently small $E$ it also lies within the fixed bracket
neighborhood of Lemma~\ref{lem:rare-coarse-normals}. Add the exact
dyadic compass shift after rounding the target once. Thus the preceding
radial conversion and normal certification apply to every physical
command. Certified conservative integer batch sizes preserve the
probability and cost bounds. The finest coordinates and all batch
descriptions require only
$O(\log(C/\nu)+\log(C/\delta))$ digits.

Finally, the radial errors are $O(h^s)$ on the common successful event.
The interpolation, support subtraction, convexity and period-locking
steps of Theorem~\ref{thm:rare-stationary} apply without change. At
two known scales, include all searches in $M_h$, use the same uniform
constants, and apply its complete-component correspondence and common
footprint reconstruction. This gives the stated output conclusions
with total failure probability at most $\delta$.
\end{proof}
